% Layout sketch for EuCAIF poster — structure/geometry only.
% Two real figures already dropped in this folder are wired in below;
% everything else is a labelled placeholder box (see \placeholderfig).
\documentclass[final]{beamer}
\usepackage[size=a0,scale=1.15,orientation=portrait]{beamerposter}
\usetheme{default}
\usepackage{tikz}
\usetikzlibrary{shadows}
\usepackage{graphicx}
\usepackage{xcolor}
\usepackage{ragged2e}

\definecolor{accent}{HTML}{1B6CA8}
\definecolor{darktext}{HTML}{222222}
\definecolor{lightbg}{HTML}{F4F7FA}
\definecolor{highlight}{HTML}{C0392B}
\definecolor{cardline}{HTML}{D9DEE3}
\definecolor{futurebg}{HTML}{EAF3FB}

\setbeamercolor{normal text}{fg=darktext,bg=white}
\setbeamercolor{block body}{fg=darktext,bg=white}
\setbeamercolor{title}{fg=accent}
% Title now lives inside the white card (see block begin/end below), set in
% the accent colour with real weight, instead of white-on-a-solid-bar.
\setbeamerfont{block title}{series=\bfseries,size=\Large}
\setbeamerfont{title}{series=\bfseries,size=\Huge}
\setbeamertemplate{itemize items}{\textcolor{accent}{\textbullet}}
\setbeamertemplate{navigation symbols}{}

% ---------------------------------------------------------------------
% "Modern card" look: one rounded, white, drop-shadowed box per block,
% title set in the accent colour with generous breathing room and a thin
% rule underneath -- no more solid-colour title bar. \cardbox{} below
% exposes the same styling standalone, for the placeholder/highlight
% boxes, so every box on the poster shares one visual language.
% ---------------------------------------------------------------------
\newlength{\cardpad}
\setlength{\cardpad}{5mm}
\newsavebox{\cardreg}
\newcommand{\cardbox}[2][cardline]{%
  \sbox{\cardreg}{#2}%
  \begin{tikzpicture}
    \node[rounded corners=7pt,draw=#1,line width=0.7pt,fill=white,
          inner sep=\cardpad,
          drop shadow={shadow xshift=0.9mm,shadow yshift=-0.9mm}]
      {\usebox{\cardreg}};
  \end{tikzpicture}%
}

\newsavebox{\blockbox}
\setbeamertemplate{block begin}{%
  \par\vskip3mm
  \begin{lrbox}{\blockbox}
  \begin{minipage}[t]{\dimexpr\linewidth-2\cardpad\relax}
  {\usebeamerfont{block title}\color{accent}\insertblocktitle}\par
  \vspace{2mm}
  {\color{accent!35}\rule{\linewidth}{0.9pt}}\par
  \vspace{4mm}
  \justifying
}
\setbeamertemplate{block end}{%
  \end{minipage}
  \end{lrbox}
  \begin{tikzpicture}
    \node[rounded corners=7pt,draw=cardline,line width=0.7pt,fill=white,
          inner sep=\cardpad,
          drop shadow={shadow xshift=0.9mm,shadow yshift=-0.9mm}]
      {\usebox{\blockbox}};
  \end{tikzpicture}
  \vskip3mm
}

% Light-tinted variant of \cardbox, for the "Future plans" callout -- same
% rounded/shadowed card language but a filled accent-blue background so it
% visually stands apart from the plain-white content blocks.
\newsavebox{\futurereg}
\newcommand{\futurebox}[1]{%
  \sbox{\futurereg}{#1}%
  \begin{tikzpicture}
    \node[rounded corners=7pt,draw=accent,line width=0.9pt,fill=futurebg,
          inner sep=\cardpad,
          drop shadow={shadow xshift=0.9mm,shadow yshift=-0.9mm}]
      {\usebox{\futurereg}};
  \end{tikzpicture}%
}

% Generic "figure goes here" box: fixed height #1 (default 6cm), full column
% width, centered italic label #2. Swap for \includegraphics once a figure
% is chosen — no other change needed.
\newcommand{\placeholderfig}[2][6cm]{%
  \cardbox{%
    \begin{minipage}[c][#1][c]{\dimexpr\linewidth-2\cardpad\relax}
      \centering\itshape\color{black!45} #2
    \end{minipage}%
  }%
}

\title{Simulation-Based inference for massive black\\[10pt] hole binary from mock LISA data}
\author{\textbf{G. Puleo}$^\dagger$, R. Trotta, E. Barausse}
\institute{Scuola Internazionale Superiore di Studi Avanzati (SISSA), Trieste, Italy}

\begin{document}
\begin{frame}[t]

% ---------------- Header ----------------
\begin{columns}[T]
\begin{column}{0.16\paperwidth}
  \centering
  \includegraphics[height=8.7cm]{sissa_logo.pdf}\hspace{0.7cm}%
  \begin{minipage}[t]{4.5cm}
    \centering
    \includegraphics[height=8.7cm]{my_profile_picture.jpg}\\[0.15cm]
    {\footnotesize $^\dagger$\texttt{gpuleo@sissa.it}}
  \end{minipage}
\end{column}
\begin{column}{0.66\paperwidth}
  \centering
  {\usebeamerfont{title}\usebeamercolor[fg]{title}\inserttitle}\\[0.6cm]
  {\Large \insertauthor}\\[0.3cm]
  {\large \insertinstitute}
\end{column}
\begin{column}{0.16\paperwidth}
  \centering
  \includegraphics[height=8.7cm]{eucaif_logo.png}
\end{column}
\end{columns}

\vspace{0.75cm}

% ---------------- Body: two columns ----------------
\begin{columns}[t]

% ===== Column 1: Motivation + Method =====
\begin{column}{0.47\paperwidth}

  \begin{block}{Motivation and outcome}
    \begin{itemize}
      \justifying
      \item The Laser Interferometer Space Antenna (LISA) will be sensitive to gravitational waves (GW) from a plethora of sources, including the coalescence of massive black hole binaries (MBHBs). 
      \item Several mergers of MBHBs are expected to be observed by LISA during operations with signal to noise ratio (SNR) of thousands. Classical inference requires a tractable likelihood; while one is available under simplified Gaussian stationary noise, more realistic assumptions -- overlapping sources, non-stationary noise, instrumental glitches -- can make it intractable or prohibitively expensive to evaluate.
      \item We present an improved version of truncated marginal neural ratio estimation (TMNRE), a simulation-based (likelihood-free) inference method, adapted to LISA-like data. We benchmark it against Markov Chain Monte Carlo (MCMC) on a signal with SNR$=2200$, produced by a non-rotating source of chirp mass $\mathcal M_{\rm c}=10^{5.25} M_{\odot}$.
      \item We successfully recover the 1D marginal posterior on the 11 source parameters, and the 2D marginal posterior on the sky location of the event. 
    \end{itemize}
  \end{block}

  \begin{block}{Observed data: LISA time series $\to$ frequency domain}
    \centering
    \includegraphics[width=0.72\linewidth]{lisa-time-freq-data-example.pdf}

    \justifying
    The LISA data is a time series in three channels (A, E, T); we train on mock A and E data, generated directly in the Fourier domain with the IMRPhenomHM waveform model (channel A shown here).
    The right panel is the frequency-domain observation used for inference; the left panel is its time-domain counterpart, obtained by inverse Fourier transform -- akin to what LISA would observe.
  \end{block}

  \begin{block}{TMNRE method: architecture}
    \centering
    \includegraphics{architecture_chmlp_tmnre_simple_crop.pdf}

  \justifying
  A neural network receives the observed data and a candidate set of source parameters, and is trained as a binary classifier: does this data actually come from these parameters, or from unrelated ones? The optimal classifier for this task can be shown to recover exactly the posterior-to-prior ratio, so evaluating the trained network on a grid of parameter values, for a fixed observation, yields the posterior.
  Here $(\lambda,\beta)$ (ecliptic longitude/latitude, i.e.\ sky location) illustrate one such marginal; in practice we train one classifier head per parameter (or pair), covering all 11 source parameters.
  \end{block}

  \begin{block}{TMNRE workflow}
    \centering
    \includegraphics[width=0.82\linewidth]{tmnre-workflow-final.pdf}

    \justifying
    Training proceeds in rounds: within a round, the simulator continuously overwrites the training set with fresh simulations, increasing the effective dataset size at a fixed memory cost. Each round ends with \emph{truncation}: the current posterior estimate defines a narrower proposal prior for the next round, focusing subsequent simulations on the region actually compatible with the data rather than the full prior volume.
  \end{block}
\end{column}

% ===== Column 2: Results =====
\begin{column}{0.47\paperwidth}

  \begin{block}{Multimodal truncation}
    \noindent
    \begin{minipage}[c]{0.542\linewidth}
      \centering
      \includegraphics[width=\linewidth]{cherry_posterior_inc_round25.pdf}
    \end{minipage}%
    \hfill
    \begin{minipage}[c]{0.418\linewidth}
      \justifying
      Our truncation algorithm accommodates multimodal posteriors, in 1D and 2D: a single bounding box would either miss a mode or waste simulations on the empty region between modes, so we track each mode's support individually. Shown: the $\cos\iota$ (inclination) and sky-location proposal posteriors and training prior at round 25.
      Angular periodicity is enforced -- note mode 1 of the sky posterior has matching density at $\lambda=0$ and $\lambda=2\pi$.
    \end{minipage}

    \vspace{1.5cm}
    \centering
    \includegraphics{cherry_posterior_sky_round25.pdf}
  \end{block}

  \begin{block}{Truncation as a function of training round}
    \centering
    \includegraphics[width=0.95\linewidth]{interval_evolution_truncation_params_logMchirp_Deltat_chi_eff.pdf}
    \par
    \vspace{6mm}
    \noindent
    \begin{minipage}[c]{0.473\linewidth}
      \centering
      \includegraphics[width=\linewidth]{volume_from_masks_global.pdf}
    \end{minipage}%
    \hfill
    \begin{minipage}[c]{0.487\linewidth}
      \justifying
      As rounds proceed, the proposal prior shrinks and stabilizes around the final posterior. Top: the prior bounds per round for a selection of three parameters, the time of coalescence $\Delta t$, the effective spin $\chi_{\rm eff}$ and the chirp mass. Bottom: the cumulative prior volume (product of all per-marginal volumes, normalised to volume of the initial prior) -- a $\sim10^{24}\times$ reduction by round 42, reached in under 30 GPU-hours using a single NVIDIA A100 GPU.
    \end{minipage}
  \end{block}

  \begin{block}{\textcolor{highlight}{Final posterior vs.\ MCMC}}
    \cardbox[highlight]{%
      \begin{minipage}{\dimexpr\linewidth-2\cardpad\relax}
        \centering
        \includegraphics[width=\linewidth]{round_42_last_posterior_vs_mcmc_truncation.pdf}

        \justifying
        \noindent
        \raisebox{-0.5cm}{%
        \begin{minipage}[c]{0.48\linewidth}
          \centering
          \includegraphics[width=13.902in]{final_posterior_zoom.pdf}
        \end{minipage}%
        }
        \hfill
        \raisebox{0.5cm}{%
        \begin{minipage}[c]{0.48\linewidth}
          \justifying
          We train our model to produce one-dimensional marginal posteriors for the 11 source parameters.
          The final-round posterior agrees with that of an MCMC sampler initialised close to the true value -- our method recovers the likelihood-based result without ever evaluating the likelihood, and starting from a broad, uninformative prior for all parameters (unfeasible for MCMC).

        \end{minipage}%
        }
      \end{minipage}%
    }
  \par
  \vspace{2cm}
  \futurebox{%
    \begin{minipage}{\dimexpr\linewidth-2\cardpad\relax}
      \textbf{\color{accent}Future plans and potential developments}
      \begin{itemize}
        \justifying
        \item Increase the observation duration at every round, in order to try an online, pre-merger parameter estimation.
        \item Add the foreground due to galactic binaries into the power spectrum.
        \item Use this method to replace importance resampling for amortized SBI pipelines.
      \end{itemize}
    \end{minipage}%
  }
  \end{block}

\end{column}

\end{columns}

\end{frame}
\end{document}
